\chapter{Récolter la prime de risque de variance}\label{ch:s2:harvesting-the-variance-risk-premium}

À la fin de 84~\% des mois depuis 1990, l'indice VIX, prix de marché de la volatilité du S\&P 500 du mois suivant, était supérieur à la volatilité qui
a suivi~: en moyenne 19.5 contre 15.4. Vendre ce prix et livrer la volatilité réalisée a été l'une des transactions les plus fiables de la finance.
Ses pertes surviennent dans l'autre sens, d'un seul coup. À la fin de février 2020, le VIX s'établissait à 40.1~; la volatilité réalisée du mois suivant était de 94.5. Sur le
marché synthétique du chapitre, un portefeuille vendeur de variance rapporte un ratio de Sharpe de 1.57 jusqu'à son pire mois, qui coûte alors 7.7 de ses écarts-types
mensuels~: 22~\% du capital pour une cible de volatilité de 10~\%, et plus que le capital à 60~\%. Ce chapitre porte sur les mises en œuvre des \omterm{def:s2:harvesting-the-variance-risk-premium:short}{stratégies vendeuses
de volatilité} et sur leur dimensionnement pour le mois qui compte. La construction repose sur deux composants~: \texttt{firm.synthvol}, le marché d'options synthétique sur lequel
tourne la partie~I, et \texttt{firm.shortvol}.

\section{La prime et les raisons de son existence}

La prime de risque de variance (livre~5, chapitre~14) est la différence entre le prix de la variance, le taux d'un swap de variance ou le carré d'un indice
comme le VIX, et la variance réalisée ensuite. Carr et Wu l'ont mesurée directement~: ils ont synthétisé les taux de swaps de variance à partir des prix d'options sur
cinq indices boursiers et 35 actions, et a retenu la différence avec la variance réalisée comme prime. Elle est versée parce que la variance est la plus élevée quand
les investisseurs sont les plus démunis~: une position longue en variance rapporte lors des krachs, et les investisseurs paient pour la détenir, comme ils paient pour une assurance.

\begin{definition}[Stratégie vendeuse de volatilité]\label{def:s2:harvesting-the-variance-risk-premium:short}
Une \emph{stratégie vendeuse de volatilité}\index{stratégie vendeuse de volatilité} perçoit la prime de risque de variance en vendant des options ou de la variance et en livrant
la volatilité réalisée~: swaps de variance vendus, options vendues couvertes en delta, calls couverts, puts vendus, ou contrats à terme vendeurs sur volatilité. Ses rendements sont faibles et
fréquents sur les marchés calmes, et importants et négatifs lorsque la volatilité bondit.
\end{definition}

Le bilan réel, d'après les données de Cboe (\cref{fig:s2:harvesting-the-variance-risk-premium:gap})~: sur 440 fins de mois de janvier 1990 à août 2026, le VIX
affichait en moyenne 19.5~\% et la volatilité réalisée du S\&P 500 du mois suivant 15.4~\%~; le VIX était supérieur dans 84~\% des cas, de 4.1 points en moyenne~; son
carré dépassait la variance réalisée de 31~\% en moyenne. Les six mois où il est resté le plus en deçà racontent l'autre face~: février 2020 (VIX 40.1, volatilité
réalisée 94.5), septembre 2008 (39.4 et 82.3), août 2008 (20.6 et 54.3), mars 2025 (22.3 et 47.7), juillet 2011 et juin 2002.

\begin{omfigure}
\begin{tikzpicture}
\begin{axis}[width=11cm,height=5.2cm,xlabel={\small année},ylabel={\small VIX moins volatilité réalisée (points)},tick label style={font=\footnotesize},
  xmin=1990,xmax=2027,ymin=-15,ymax=15,grid=major,grid style={gray!25},table/col sep=comma,/pgf/number format/1000 sep={}]
\addplot[omThm,thick] table[x=year,y=gap_pts]{figdata/strategies-2/01-harvesting-the-variance-risk-premium/gap.csv};
\addplot[black!40,dashed] coordinates {(1990,0) (2027,0)};
\end{axis}
\end{tikzpicture}

\omcaption{La prime de variance du S\&P 500, 1990--2026~: moyenne sur douze mois, à chaque fin de mois, du VIX moins la volatilité réalisée du S\&P
500 sur les 21 séances suivantes. Dérivée des historiques de Cboe pour le VIX et le S\&P 500~; les séries brutes ne sont pas redistribuées. Données~:
\texttt{s2\_fetch\_cboe}.}\label{fig:s2:harvesting-the-variance-risk-premium:gap}
\end{omfigure}

\section{Mises en œuvre~: variance, options, vente d'options couvertes}

\texttt{firm.synthvol} (\cref{lst:s2:harvesting-the-variance-risk-premium:synthvol}) simule vingt ans d'un indice et de trente valeurs qui le composent~: un facteur commun
à variance stochastique autour de 16~\% (retour à la moyenne de quatre par an, corrélation de levier de $-0.7$), des sauts négatifs une fois par an en moyenne, et
trois krachs implantés de 25~\%, 20~\% et 30~\% sur dix jours chacun. La loi sous-jacente fixe la surface~: la variance implicite à la monnaie de l'indice est l'espérance
de la variance sur la durée de vie de l'option, multipliée par 1.3, et le smile de volatilité décroît avec le strike. La volatilité réalisée sur les vingt ans est de 19.0~\%, krachs inclus, et
la volatilité implicite moyenne à un mois s'établit à 18.3~\%.

\begin{definition}[Vente d'options couvertes]\label{def:s2:harvesting-the-variance-risk-premium:overwrite}
La \emph{vente d'options couvertes}\index{vente d'options couvertes} (overwriting) consiste à vendre des options contre une position déjà détenue, le plus souvent des calls contre une position longue en actions
(calls couverts), en encaissant une prime en échange de la renonciation aux gains au-delà du strike~; c'est une \omterm{def:s2:harvesting-the-variance-risk-premium:short}{stratégie vendeuse de volatilité} à laquelle s'ajoute la direction
du marché.
\end{definition}

\texttt{firm.shortvol} fait tourner cinq mises en œuvre mensuelles sur l'indice, chacune renouvelée toutes les 21 séances de cotation~:

\begin{center}
\small
\begin{tabular}{@{}p{5.2cm}>{\centering\arraybackslash}p{2.2cm}>{\centering\arraybackslash}p{2.2cm}>{\centering\arraybackslash}p{2.6cm}>{\centering\arraybackslash}p{2.2cm}@{}}
\toprule
mise en œuvre, 239 mois & ratio de Sharpe & asymétrie & pire mois (écart-type) & mois en hausse\\
\midrule
swap de variance vendu & 0.69 & $-3.53$ & $-7.7$ & 79~\%\\
straddle vendu couvert en delta & 1.24 & $-2.32$ & $-6.4$ & 77~\%\\
calls couverts, 2~\% hors de la monnaie & 0.40 & $-4.24$ & $-8.1$ & 69~\%\\
puts vendus, 5~\% hors de la monnaie & 0.08 & $-7.22$ & $-10.2$ & 92~\%\\
variance vendue dimensionnée par $1/\text{prix d'exercice}$ & 0.09 & $-12.45$ & $-14.2$ & 79~\%\\
l'indice lui-même & 0.26 & $-1.82$ & $-6.8$ & 62~\%\\
\bottomrule
\end{tabular}
\end{center}

Chaque mise en œuvre présente une asymétrie négative, et cette forme correspond à la source de la prime~: le put vendu gagne dans 92~\% des mois et perd dix
écarts-types en un seul mois. Dimensionner la variance vendue à l'inverse de son strike (davantage quand la volatilité est bon marché) la rend pire~: les krachs introduits partent
de marchés calmes, quand le portefeuille est le plus large. Les krachs synthétiques sont plus sévères que l'histoire récente~: sur les indices de Cboe de février 2007 à septembre 2026,
PutWrite a rapporté 7.2~\% par an à 13.8~\% de volatilité (un ratio de Sharpe de 0.57 sans taux sans risque) et BuyWrite 6.0~\% à 14.3~\% (0.48), contre
9.0~\% à 19.7~\% (0.53) pour l'indice de cours du S\&P 500, qui omet les dividendes. PutWrite a perdu 27.6~\% de septembre à novembre 2008 et 19.8~\% en février
et mars 2020~; BuyWrite 28.5~\% et 21.3~\%.

\section{Rendements d'options couvertes en delta}

\begin{definition}[Rendement d'option couverte en delta]\label{def:s2:harvesting-the-variance-risk-premium:hedged}
Un \emph{rendement d'option couverte en delta}\index{rendement d'option couverte en delta} est le gain d'une position d'options dont le delta est couvert dans le sous-jacent à
intervalles réguliers, rapporté à sa prime~; dans un modèle sans prime de risque de volatilité, son espérance est proche de zéro, de sorte que son signe et sa taille mesurent
la prime.
\end{definition}

Bakshi et Kapadia ont exactement utilisé cela~: sur les options sur le S\&P 500, une option longue couverte en delta a rapporté moins que zéro, un peu moins hors de la monnaie, davantage quand
la volatilité était élevée, même après prise en compte des craintes de sauts. Côté vendeur, un straddle vendu à la monnaie, couvert chaque jour à la volatilité
courante de la surface (\cref{lst:s2:harvesting-the-variance-risk-premium:shortvol}), rapporte 0.53~\% de l'indice par mois sur le marché synthétique, soit environ un huitième de la
prime du straddle, de 4.2~\%, avec un ratio de Sharpe de 1.24, le meilleur des cinq, parce que la couverture supprime la direction propre de l'indice et laisse l'écart entre
la volatilité implicite et la volatilité réalisée.

\section{Effondrements et dimensionnement pour y faire face}

Vendre de la volatilité est un problème de dimensionnement. La \Cref{fig:s2:harvesting-the-variance-risk-premium:paths} met à l'échelle le portefeuille vendeur de variance pour trois cibles de
volatilité. Son ratio de Sharpe avant le pire mois (le premier krach, à l'année 5.9) vaut 1.57 quelle que soit la taille~; le pire mois coûte 7.7 de ses écarts-types
mensuels quelle que soit la taille~:

\begin{center}
\small
\begin{tabular}{@{}lccccc@{}}
\toprule
volatilité cible & 5~\% & 10~\% & 20~\% & 40~\% & 60~\%\\
\midrule
ratio de Sharpe avant le pire mois & 1.57 & 1.57 & 1.57 & 1.57 & 1.57\\
pire mois, part du capital & $-11.1~\%$ & $-22.2~\%$ & $-44.3~\%$ & $-88.6~\%$ & $-132.9~\%$\\
le capital survit & oui & oui & oui & oui & non\\
\bottomrule
\end{tabular}
\end{center}

Le passé qui précède le krach ne dit rien de sa taille~: un gérant qui jugeait le risque à partir de cinq années de rendements mensuels voyait un ratio de Sharpe
supérieur à 1.5 et un pire mois d'une fraction de celui à venir. La taille qui survit est fixée par le scénario de stress, et non par la volatilité~: un portefeuille qui ne doit pas
perdre plus d'un quart de son capital en un mois peut faire tourner le portefeuille synthétique vendeur de variance à une volatilité d'environ 10~\%, quel que soit son historique.

\begin{omfigure}
\begin{tikzpicture}
\begin{axis}[width=11cm,height=5.4cm,ymode=log,xlabel={\small année},ylabel={\small capital (échelle logarithmique)},tick label style={font=\footnotesize},
  xmin=0,xmax=20,ymin=0.3,ymax=50,ytick={0.5,1,2,5,10,20,50},yticklabels={0.5,1,2,5,10,20,50},grid=major,grid style={gray!25},legend columns=3,legend style={font=\scriptsize,at={(0.5,-0.3)},anchor=north,
  draw=none,fill=none,/tikz/every even column/.append style={column sep=8pt}},table/col sep=comma]
\addplot[omExo,thick] table[x=year,y=v10]{figdata/strategies-2/01-harvesting-the-variance-risk-premium/paths.csv};
\addplot[omDef,thick] table[x=year,y=v20]{figdata/strategies-2/01-harvesting-the-variance-risk-premium/paths.csv};
\addplot[omThm,thick] table[x=year,y=v40]{figdata/strategies-2/01-harvesting-the-variance-risk-premium/paths.csv};
\legend{cible 10~\%,cible 20~\%,cible 40~\%}
\end{axis}
\end{tikzpicture}

\omcaption{Le portefeuille synthétique vendeur de variance, dimensionné à trois cibles de volatilité et capitalisé mensuellement sur vingt ans~; les krachs programmés des années 6.0 et
17.5 apparaissent comme des baisses en un seul mois, le plus petit, celui de l'année 12.7, à peine. Données~: \texttt{s2\_vrp.paths}.}\label{fig:s2:harvesting-the-variance-risk-premium:paths}
\end{omfigure}

\section{Fiches de stratégie}

\begin{strategyfile}{Vente de variance indicielle}\label{strat:s2:harvesting-the-variance-risk-premium:variance}
\sfield{Qui paie, et pourquoi} Les investisseurs qui paient pour se protéger contre la variance, laquelle est la plus élevée dans les mauvais moments.
\sfield{Instruments et plateformes} Swaps de variance de gré à gré~; un panier d'options cotées.
\sfield{Signal} La prime de variance~: la variance implicite comparée à une prévision de la variance réalisée.
\sfield{Dimensionnement et exécution} Dimensionnée par un stress (le pire mois plausible), et non par la volatilité~; roll mensuel.
\sfield{Coûts} Fourchettes des options du panier de réplication~; marges des dealers sur les swaps.
\sfield{Comment elle meurt} Les krachs~; l'encombrement (crowding), qui accroît la taille du krach.
\sfield{Horizon, capacité, infrastructure} Un mois~; une grande capacité sur la variance indicielle~; une surface de volatilité et un système de stress.
\sfield{Backtest honnête} Les années de krach dans l'échantillon~; les strikes de variance tels qu'ils sont cotés~; les appels de marge.
\sfield{Sources} Carr et Wu (2009)~; données VIX du Cboe~; ce chapitre~: 1.57 avant le pire mois, $-7.7$ écarts types dans ce mois.
\end{strategyfile}

\begin{strategyfile}{Straddle vendeur couvert en delta}\label{strat:s2:harvesting-the-variance-risk-premium:straddle}
\sfield{Qui paie, et pourquoi} Comme pour la vente de variance, avec la direction couverte.
\sfield{Instruments et plateformes} Options indicielles à la monnaie cotées et contrats à terme pour la couverture.
\sfield{Signal} Volatilité implicite comparée à la volatilité prévue.
\sfield{Dimensionnement et exécution} Vendu chaque mois~; couvert en delta chaque jour ou sur des bandes.
\sfield{Coûts} Fourchettes des options et négociation de la couverture.
\sfield{Comment elle meurt} Les sauts de cours que la couverture ne peut pas suivre.
\sfield{Horizon, capacité, infrastructure} Un mois~; un moteur de couverture.
\sfield{Backtest honnête} Couvertures aux prix exécutables~; sauts de cours et suspensions de cotation.
\sfield{Sources} Bakshi et Kapadia (2003)~; ce chapitre~: un ratio de Sharpe de 1.24 sur l'indice synthétique.
\end{strategyfile}

\begin{strategyfile}{Vente d'options couvertes}\label{strat:s2:harvesting-the-variance-risk-premium:overwrite}
\sfield{Qui paie, et pourquoi} Les acheteurs de calls sur la hausse~; la prime de variance sur les calls.
\sfield{Instruments et plateformes} Une position en actions et des calls cotés sur ces actions.
\sfield{Signal} Aucun, hormis la prime~; strike et échéance choisis selon une règle.
\sfield{Dimensionnement et exécution} Calls vendus chaque mois sur la totalité de la position.
\sfield{Coûts} Fourchettes des options~; la hausse abandonnée.
\sfield{Comment elle meurt} Elle ne meurt pas~; elle sous-performe dans les fortes hausses et baisse avec le marché.
\sfield{Horizon, capacité, infrastructure} Un mois~; une très grande capacité.
\sfield{Backtest honnête} Inclure les dividendes dans l'indice de comparaison.
\sfield{Sources} Indice Cboe BuyWrite~: 6.0~\% par an à une volatilité de 14.3~\%, 2007--2026.
\end{strategyfile}

\begin{strategyfile}{Vente de puts}\label{strat:s2:harvesting-the-variance-risk-premium:putwrite}
\sfield{Qui paie, et pourquoi} Les acheteurs de protection contre les krachs, qui paient le plus cher pour les puts hors de la monnaie.
\sfield{Instruments et plateformes} Puts indiciels cotés, collatéralisés en espèces.
\sfield{Signal} Aucun, hormis la prime.
\sfield{Dimensionnement et exécution} Entièrement collatéralisée~; puts à un mois roulés à l'échéance.
\sfield{Coûts} Fourchettes des options.
\sfield{Comment elle meurt} Les krachs~: le put vendu gagne la plupart des mois et perd sur quelques-uns.
\sfield{Horizon, capacité, infrastructure} Un mois~; une très grande capacité.
\sfield{Backtest honnête} Collatéral rémunéré au taux des bons du Trésor~; années de krach dans l'échantillon.
\sfield{Sources} Indice Cboe PutWrite~: 7.2~\% par an à une volatilité de 13.8~\%, 2007--2026~; les krachs synthétiques plus sévères de ce chapitre.
\end{strategyfile}

\begin{strategyfile}{Vente de volatilité ciblée}\label{strat:s2:harvesting-the-variance-risk-premium:voltarget}
\sfield{Qui paie, et pourquoi} Comme pour la vente de variance.
\sfield{Instruments et plateformes} Swaps de variance ou straddles.
\sfield{Signal} La prime, dimensionnée par la volatilité implicite courante.
\sfield{Dimensionnement et exécution} Plus grande lorsque la volatilité implicite est basse, plus petite lorsqu'elle est haute.
\sfield{Coûts} Rotation supplémentaire du portefeuille.
\sfield{Comment elle meurt} Les krachs qui partent de marchés calmes, lorsque la position est la plus grande.
\sfield{Horizon, capacité, infrastructure} Un mois.
\sfield{Backtest honnête} Krachs partant de départs à faible volatilité dans l'échantillon.
\sfield{Sources} Ce chapitre~: un dimensionnement par l'inverse du strike a ramené le ratio de Sharpe de 0.69 à 0.09.
\end{strategyfile}

\section{Tutoriel~: cueillir des centimes}\label{tut:s2:harvesting-the-variance-risk-premium:tut}

\begin{tutorial}
\textbf{Objectif.} Construire le marché synthétique d'options, exécuter cinq implémentations de vente de volatilité et dimensionner le portefeuille vendeur de variance à travers ses
krachs programmés~; comparer avec la prime de variance réelle et les indices d'écriture d'options. \textbf{État final~:} les tableaux et les deux figures.
\begin{tutsteps}
\item \textbf{Le marché}~: variance stochastique avec effet de levier, sauts et krachs~; composantes à rendements spécifiques.
\omcode{code/firm/synthvol/firm_synthvol.py}{95}{120}{Vingt ans de l'indice synthétique et de ses composantes.}\label{lst:s2:harvesting-the-variance-risk-premium:synthvol}
\item \textbf{Deux implémentations}~: le swap de variance et le straddle couvert en delta.
\omcode{code/firm/shortvol/firm_shortvol.py}{33}{58}{Vente de variance et straddle vendeur couvert en delta.}\label{lst:s2:harvesting-the-variance-risk-premium:shortvol}
\item \textbf{Exécuter} \texttt{s2\_fetch\_cboe.py} une fois pour les statistiques réelles, puis \texttt{implementations()}, \texttt{leverage\_table()} et
\texttt{fig\_vrp.py}.
\end{tutsteps}
\textbf{Pour aller plus loin.} Couvrir le straddle sur bandes plutôt que chaque jour~; ajouter un krach qui part d'une volatilité élevée et vérifier si le dimensionnement
par l'inverse du strike aide~; acheter un put très hors de la monnaie contre la variance vendue et mesurer ce qu'il coûte et ce qu'il évite.
\end{tutorial}

\section{Construction~: le marché synthétique d'options et la vente de volatilité}\label{bld:s2:harvesting-the-variance-risk-premium:shortvol}

\begin{build}
\bfield{Rôle} \texttt{firm.synthvol}~: un marché d'options déterministe sur indice et sur composantes, avec variance stochastique, sauts, krachs et surfaces portant
une prime de variance et une prime de corrélation, pour la partie~I. \texttt{firm.shortvol}~: implémentations de vente de volatilité et effet de levier.
\bfield{Interface} \texttt{VolConfig(\dots)}, \texttt{simulate\_vol(cfg)}, \texttt{expected\_var}, \texttt{atm\_iv}, \texttt{smile\_iv}, \texttt{bs\_price},
\texttt{bs\_delta}~; \texttt{periods}, \texttt{short\_variance}, \texttt{hedged\_straddle}, \texttt{overwrite}, \texttt{put\_write}, \texttt{lever}.
\bfield{Règles} Volatilités implicites tirées de la vérité connue à l'époque~; options évaluées sur la surface~; P\&L par unité d'indice à l'initiation.
\bfield{Tests d'acceptation} \texttt{code/firm/synthvol/tests/} et \texttt{code/firm/shortvol/tests/}~: Black--Scholes à la main, variance attendue et
prime, un krach programmé~; P\&L de la variance et du straddle sur un chemin plat, une ruine totale.
\bfield{Pour aller plus loin} Surfaces avec primes de structure par terme~; coûts de transaction sur les couvertures~; options des composantes.
\end{build}

\begin{omsources}
\item P.~Carr et L.~Wu, «~Variance risk premiums~», \emph{Review of Financial Studies} 22(3), 2009.
\item G.~Bakshi et N.~Kapadia, «~Delta-hedged gains and the negative market volatility risk premium~», \emph{Review of Financial Studies} 16(2), 2003.
\item Cboe Global Markets, historiques quotidiens du VIX, du S\&P 500, des indices PutWrite et BuyWrite.
\end{omsources}

\section{Exercices}

\begin{exercise}[$\star$]\label{exo:s2:harvesting-the-variance-risk-premium:1}
Un straddle à la monnaie vaut environ $0.8\,\sigma\sqrt{\tau}$ de l'indice. Calculer sa valeur pour un mois (21 jours de négociation) à 20~\% de volatilité.
\end{exercise}

\begin{exercise}[$\star$]\label{exo:s2:harvesting-the-variance-risk-premium:2}
La variance attendue vaut $0.16^2$ et la variance implicite est supérieure de 30~\% à celle-ci. Calculer la variance implicite et la volatilité implicite.
\end{exercise}

\begin{exercise}[$\star$]\label{exo:s2:harvesting-the-variance-risk-premium:3}
Expliquer pourquoi le put vendu gagne plus souvent que le swap de variance vendu, mais affiche un ratio de Sharpe plus faible.
\end{exercise}

\begin{exercise}[$\star\star$]\label{exo:s2:harvesting-the-variance-risk-premium:4}
Un portefeuille à 10~\% de volatilité annuelle perd 7.7 écarts-types mensuels. Calculer la fraction du capital qu'il perd.
\end{exercise}

\begin{exercise}[$\star\star$]\label{exo:s2:harvesting-the-variance-risk-premium:5}
Expliquer pourquoi dimensionner la vente de variance par l'inverse de son strike dégrade le résultat sur le marché synthétique. Dans quels cas cette méthode pourrait-elle aider~?
\end{exercise}

\begin{exercise}[$\star\star$]\label{exo:s2:harvesting-the-variance-risk-premium:6}
Expliquer pourquoi le VIX dépasse la volatilité réalisée la plupart du temps, et pourquoi son excès moyen ne mesure pas à lui seul la totalité de la prime.
\end{exercise}

\begin{exercise}[$\star\star\star$]\label{exo:s2:harvesting-the-variance-risk-premium:7}
\emph{Programmation.} Exécuter \texttt{leverage\_table((0.3, 0.5))}. Déterminer la volatilité cible comprise entre ces deux bornes pour laquelle le pire mois emporte la totalité du capital.
\end{exercise}

\begin{exercise}[$\star\star\star$]\label{exo:s2:harvesting-the-variance-risk-premium:8}
\emph{Trouver la faille.} «~Notre fonds vendeur de volatilité affiche un ratio de Sharpe de 2 sur cinq ans, et son pire mois a été de $-3~\%$~; à 15~\% de volatilité, le risque
est modeste.~»
\end{exercise}

\section{Problème~: Ramasser les petits gains}

\begin{problem}[{Problème du week-end --- stratégie vendeuse de volatilité, dimensionnée}]\label{pb:s2:harvesting-the-variance-risk-premium:1}
Le marché d'options synthétique du chapitre, les données de Cboe et les archives publiques.

\medskip\noindent\textbf{Partie I --- La prime.}
\begin{enumerate}
\item Définir une \omterm{def:s2:harvesting-the-variance-risk-premium:short}{stratégie vendeuse de volatilité} et la \omterm{def:s2:harvesting-the-variance-risk-premium:overwrite}{vente d'options couvertes} (overwriting).
\item Que mesurent Carr et Wu, et comment~?
\item Donner le bilan du VIX face à la volatilité réalisée, de 1990 à 2026.
\item Citer les mois où la volatilité réalisée a le plus dépassé le VIX.
\end{enumerate}

\medskip\noindent\textbf{Partie II --- Implémentations.}
\begin{enumerate}[resume]
\item Décrire le marché synthétique et la manière dont sa surface est fixée.
\item Donner les ratios de Sharpe et l'asymétrie des cinq implémentations.
\item Pourquoi toutes les asymétries sont-elles négatives~?
\item Combien ont rapporté les indices PutWrite et BuyWrite, et comment se sont-ils comportés en 2008 et en 2020~?
\end{enumerate}

\medskip\noindent\textbf{Partie III --- Options couvertes.}
\begin{enumerate}[resume]
\item Définir un \omterm{def:s2:harvesting-the-variance-risk-premium:hedged}{rendement d'option couverte en delta}.
\item Que Bakshi et Kapadia ont-ils mis en évidence~?
\item Combien rapporte le straddle synthétique, et pourquoi son ratio de Sharpe est-il le plus élevé~?
\item Que changerait une couverture sur des bandes~?
\end{enumerate}

\medskip\noindent\textbf{Partie IV --- Le verdict.}
\begin{enumerate}[resume]
\item Énoncer le \emph{résultat nommé}~: le ratio de Sharpe du portefeuille vendeur de variance avant son pire mois et sa perte durant ce mois, selon l'effet de levier.
\item Pourquoi le bilan antérieur au krach ne mesure-t-il pas le krach~?
\item Comment faut-il dimensionner un portefeuille vendeur de volatilité~?
\item Pourquoi le dimensionnement par l'inverse du strike a-t-il échoué ici~?
\item Comment backtester honnêtement une \omterm{def:s2:harvesting-the-variance-risk-premium:short}{stratégie vendeuse de volatilité}~?
\item Quelle fiche de stratégie gagne le plus souvent, et lequel a le meilleur ratio de Sharpe~?
\item Quel lien ce chapitre entretient-il avec le ciblage de volatilité (livre~8, chapitre~26)~?
\item En une phrase~: que vend un trader en \omterm{def:s2:harvesting-the-variance-risk-premium:short}{stratégies vendeuses de volatilité}~?
\end{enumerate}
\end{problem}

\section{Questions d'entretien}

\begin{interviewq}[$\star$ \iqroles{chercheur}]\label{iq:s2:harvesting-the-variance-risk-premium:1}
Qu'est-ce que la prime de risque de variance, et pourquoi existe-t-elle~?
\end{interviewq}

\begin{interviewq}[$\star\star$ \iqroles{trader}]\label{iq:s2:harvesting-the-variance-risk-premium:2}
Vous êtes vendeur d'un straddle couvert en delta et l'indice fait un gap à la baisse de 5~\% pendant la nuit. Que s'est-il passé pour votre P\&L, et que faites-vous~?
\end{interviewq}

\begin{interviewq}[$\star\star$ \iqroles{risque}]\label{iq:s2:harvesting-the-variance-risk-premium:3}
Comment fixeriez-vous des limites sur un portefeuille vendeur de volatilité~?
\end{interviewq}

\begin{interviewq}[$\star\star$ \iqroles{chercheur}]\label{iq:s2:harvesting-the-variance-risk-premium:4}
Comparez la vente de swaps de variance, de straddles et d'options de vente hors de la monnaie comme moyens de capter la prime.
\end{interviewq}

\begin{interviewq}[$\star\star$ \iqroles{développeur}]\label{iq:s2:harvesting-the-variance-risk-premium:5}
De quoi a besoin un moteur quotidien de couverture en delta pour un portefeuille d'options~?
\end{interviewq}

\begin{interviewq}[$\star\star\star$ \iqroles{chercheur}]\label{iq:s2:harvesting-the-variance-risk-premium:6}
Une option couverte en delta, de gamma $\Gamma$ sur un sous-jacent valant $S$, fait l'objet d'une couverture continue. Montrer que son P\&L sur un court intervalle vaut environ
$\tfrac12\Gamma S^2(\sigma_r^2 - \sigma_i^2)\,dt$, où $\sigma_r$ désigne la volatilité réalisée et $\sigma_i$ la volatilité implicite, et expliquer ce que cela implique pour son signe.
\end{interviewq}
